\section{Implementación de un ReAct Agent con LangChain}

\subsection{Diseño de la plantilla de Prompt}

En LangChain, \texttt{create\_react\_agent} usa una plantilla de Prompt estandarizada, cuyo diseño enfatiza la \textbf{generalidad y la universalidad}:

\begin{quote}
\textit{Answer the following questions as best you can. You have access to the following tools: \{tools\}. Use the following format: Question / Thought / Action / Action Input / Observation / ... / Final Answer. Begin! Question: \{input\} \{agent\_scratchpad\}}
\end{quote}

La plantilla contiene dos marcadores dinámicos:
\begin{itemize}
    \item \texttt{tools}: la lista de herramientas disponibles
    \item \texttt{agent\_scratchpad}: el historial que se mantiene dinámicamente durante la interacción
\end{itemize}

\subsection{Proceso de construcción en tres pasos}

\begin{enumerate}
    \item Declarar el modelo de lenguaje (por ejemplo, GPT-4.1 Nano)
    \item Llamar a \texttt{create\_react\_agent} para crear el Agent
    \item Colocar el Agent dentro de un \texttt{AgentExecutor} para ejecutarlo
\end{enumerate}

\subsection{Mecanismo de funcionamiento: iteración basada en una plantilla de texto}

La implementación de ReAct en LangChain tiene una característica importante: \textbf{no usa una lista de mensajes de varios turnos (message list)}, sino que va agregando toda la información del historial a \textbf{un mismo mensaje human}:

\begin{warningbox}{Forma en que LangChain ReAct administra el historial}
Cada nuevo turno de interacción es una nueva llamada al LLM. La información del historial (Thought, Action, Observation) se agrega a \texttt{agent\_scratchpad} y se envía al modelo como parte de un único mensaje human. Esto significa que no hay una estructura de mensajes de varios turnos: todo el contexto queda apretado en un solo mensaje.
\end{warningbox}

\subsection{Resumen de la sección}
La implementación de ReAct en LangChain administra el historial de la interacción mediante una plantilla de Prompt más la concatenación de texto; la estructura es simple, pero carece de la organización clara de los mensajes de varios turnos.

